\begin{table}[htbp]\centering\small
\caption{Gasto: agregados del decreto y de los Criterios Generales (mdp)}\label{cua:gasto}
\begin{tabular}{lrrr}
\toprule
Concepto & 2026 & 2027 & Var. real (\%) \\
\midrule
Gasto neto total (decreto, línea P) & 10.193.683,7 & 10.636.488,1 & +0,3 \\
Neteo (decreto, línea P) & 1.553.113,1 & 1.680.127,3 & +4,0 \\
Gasto neto pagado (CGPE, línea G) & 10.114.827,9 & 10.515.087,7 & $-$0,04 \\
\quad Programable pagado (línea G) & 7.015.853,1 & 7.311.129,3 & +0,2 \\
\quad No programable (línea G) & 3.098.974,9 & 3.203.958,4 & $-$0,6 \\
\quad\quad Costo financiero (línea G) & 1.572.073,3 & 1.574.950,2 & $-$3,7 \\
Gastos obligatorios sin pensiones (línea P) & 6.702.281,0 & 6.693.540,5 & $-$4,0 \\
Perímetro de pensiones, por diferencia (línea P) & 1.704.159,2 & 1.840.746,1 & +3,9 \\
\bottomrule
\end{tabular}
\\[2pt]\begin{minipage}{0.92\textwidth}\scriptsize \textbf{Las dos líneas no se mezclan y van etiquetadas renglón por renglón.} Fuentes: Anexos 1 y 3 de los dos proyectos de decreto; CGPE 2027 Anexo II.6. El perímetro de pensiones sale de la diferencia entre los dos renglones del Anexo 3, no de una selección propia de programas.\end{minipage}
\end{table}
